\SourceLeaf{470}
\par\Needspace{9\baselineskip}\bigskip
\pdfbookmark[3]{Exposition, 20}{pallas-II-article20}
\begin{center}20.\end{center}
\begingroup
\setlength{\lineskip}{4pt}\setlength{\lineskiplimit}{4pt}
\noindent\hspace*{\parindent}The variable values of the elements obtained by this method will be accurate up to quantities of the order of the square of the perturbing masses. If still greater precision is desired, one need only repeat this calculation using the variable elements just found. This gives results accurate up to quantities of the third order. It would be superfluous to repeat the calculation once more. In the second calculation one may even assume. that the elements remain constant for some time, i.e. use one and the same set of elements in calculating the instantaneous variations for several epochs before or after the one to which it corresponds. In general, one does not change the set until the difference between it and the set that would have to replace it becomes appreciable enough. Consequently, the smaller the instantaneous variations in the elements, the greater the number of epochs for which the same set of elements may be used.
\par In practice it will be a little more convenient to multiply the values of \(\dfrac{\mathrm{d}i}{\mathrm{d}t},\ \dfrac{\mathrm{d}\AscendingNode}{\mathrm{d}t},\ \dfrac{\mathrm{d}\varphi}{\mathrm{d}t},\ \dfrac{\mathrm{d}\tilde\omega}{\mathrm{d}t},\ \dfrac{\mathrm{d}\epsilon}{\mathrm{d}t}\) by \(\theta\) and those of \(\dfrac{\mathrm{d}n}{\mathrm{d}t}\) by \(\theta\theta\) before forming the series of differences and sums. It hardly needs saying that these quantities should first be expressed in seconds. The most suitable size of the interval \(\theta\) depends partly on the speed of motion of the perturbing and perturbed planets and partly on the degree of precision sought. In using this method to calculate the perturbations of Pallas produced by Jupiter during the years \mbox{1803 . . . . 1816}, the author chose an interval of 50 days. The discerning calculator must judge whether, in other circumstances, this interval should be decreased or increased. We must omit here several other practical remarks that will suggest themselves to any experienced calculator undertaking similar work.
\par\endgroup
\begin{center}\rule{0.18\textwidth}{0.35pt}\end{center}
